% ===========================================================================
\chapter{Les sorties}
% ===========================================================================

% ─────────────────────────────────────────────────────────────────────────────
\section{Le CSV de r\'esultats}
% ─────────────────────────────────────────────────────────────────────────────

Le fichier \fichier{resultats/portees\_admissibles.csv} est la sortie principale.
Chaque ligne correspond \`a une combinaison
(produit $\times$ config\_calcul $\times$ port\'ee $\times$ combinaison EN~1990).

Colonnes principales~:

\begin{center}
\begin{tabular}{ll}
\toprule
\textbf{Colonne} & \textbf{Description}\\
\midrule
\py{id\_produit} & R\'ef\'erence du produit bois\\
\py{id\_config\_calcul} & Identifiant du sc\'enario EC5\\
\py{type\_poutre} & Panne, Solive, Chevron, Sommier\\
\py{longueur\_m} & Port\'ee v\'erifi\'ee (m)\\
\py{statut} & \py{admis} / \py{refus\'e} / \py{rejet\'e\_usage}\\
\py{taux\_determinant} & Taux de sollicitation maximal\\
\py{verification\_determinante} & Quelle v\'erification est critique\\
\py{clause\_EC5} & R\'ef\'erence normative (\S{}6.1.6, \S{}7.2, etc.)\\
\py{sigma\_m\_MPa} & Contrainte de flexion calcul\'ee\\
\py{taux\_flexion\_ELU} & Taux ELU flexion\\
\py{tau\_MPa} & Contrainte de cisaillement calcul\'ee\\
\py{taux\_cisaillement\_ELU} & Taux ELU cisaillement\\
\py{taux\_appui\_ELU} & Taux ELU appui\\
\py{longueur\_appui\_min\_mm} & Longueur d'appui minimale calcul\'ee\\
\py{k\_crit} & Coefficient de d\'eversement\\
\py{taux\_deversement\_ELU} & Taux ELU d\'eversement\\
\py{w\_inst\_mm} & Fl\`eche instantan\'ee (mm)\\
\py{w\_fin\_mm} & Fl\`eche finale avec fluage (mm)\\
\py{taux\_ELS\_inst} & Taux ELS fl\`eche instantan\'ee\\
\py{taux\_ELS\_fin} & Taux ELS fl\`eche finale\\
\bottomrule
\end{tabular}
\end{center}

\begin{maitre}
Comment le programme sait-il quelle v\'erification est la plus critique~?
\end{maitre}

\begin{apprenti}
Il prend le taux le plus \'elev\'e parmi toutes les v\'erifications~?
\end{apprenti}

\begin{maitre}
Exactement~:
\end{maitre}

\begin{lstlisting}[style=python, caption={Extrait de \fichier{moteur.py} --- vérification déterminante}]
def _trouver_verification_determinante(elu, els, df, taux_max):
    """Trouve quelle vérification a le taux le plus élevé."""
    candidates = {
        "flexion_ELU":      elu.get("taux_flexion_ELU", 0),
        "cisaillement_ELU": elu.get("taux_cisaillement_ELU", 0),
        "appui_ELU":        elu.get("taux_appui_ELU", 0),
        "deversement_ELU":  elu.get("taux_deversement_ELU", 0),
        "fleche_inst":      els.get("taux_ELS_inst", 0),
        "fleche_fin":       els.get("taux_ELS_fin", 0),
    }
    # max() sur un dict : trouve la clé avec la valeur maximale
    return max(candidates, key=lambda k: candidates[k] or 0)

# Correspondance vérification -> clause normative
mapping_clause = {
    "flexion_ELU":      "§6.1.6",
    "cisaillement_ELU": "§6.1.7",
    "appui_ELU":        "§6.1.5",
    "deversement_ELU":  "§6.3.3",
    "fleche_inst":      "§7.2",
    "fleche_fin":       "§7.2",
}
\end{lstlisting}
